% =============================================================================
%  ERIT Benchmark — Experiments and Results
%  Merged methodology + results: each subsection describes the experiment
%  design and reports findings together.
% =============================================================================

\section{Experiments and Results}
\label{sec:results}

We evaluate ERIT as a \emph{benchmark suite} along three orthogonal axes---
modality, model regime, and input robustness---with two objectives:
(1)~assess whether multimodal fusion improves emotion recognition under
controlled perturbation conditions, and
(2)~analyze how current vision--language models integrate visual and
textual cues for older adult emotion recognition.

Each sample is a fixed image--text pair; since multiple emotion labels may
apply, we treat the task as multi-label classification.  For each
evaluation cell we report:
\textbf{accuracy~(Acc)}, the fraction of predictions matching the primary
ground-truth label;
\textbf{balanced accuracy~(BAcc)}, mean per-class recall;
\textbf{macro~F1}, equally weighting all seven emotions; and
\textbf{multi-label accuracy~(MAcc)}, counting a prediction correct if it
matches \emph{any} label in the original annotation set.
MAcc is used as the primary metric throughout.
Statistical significance is assessed via McNemar's test ($\alpha = 0.05$)
with bootstrap 95\,\% confidence intervals ($B = 10{,}000$).


% ─── 4.1 Baseline Results ────────────────────────────────────────────────

\subsection{Baseline Results}
\label{sec:baseline-results}

Every model is evaluated under three input configurations:
\textbf{image-only~(I)}, receiving only the representative frame;
\textbf{text-only~(T)}, receiving only the Whisper transcript (with a dummy
gray image for architectures that require visual input); and
\textbf{fusion~(I+T)}, receiving both.  We evaluate two families of
vision--language models in two regimes:

\begin{itemize}[nosep]
  \item \textbf{Zero-shot}: LLaVA-NeXT~(Mistral-7B)~\cite{liu2024llavanext}
        and Qwen2-VL~(7B-Instruct)~\cite{wang2024qwen2vl}, prompted to
        classify emotions without task-specific training.  Both models are
        evaluated on all 2{,}379 samples.
  \item \textbf{Fine-tuned}: LLaVA-NeXT with rank-8 LoRA adapters
        ($\alpha\!=\!16$, dropout$\,=\,0.1$, targeting \texttt{q\_proj} and
        \texttt{v\_proj} in the last 4 transformer layers;
        ${\sim}$426K trainable parameters) trained on the ERIT training
        partition (1{,}169 samples) with 4-bit quantisation.  The dev set
        (634 samples) is used for checkpoint selection; final results are
        reported on the held-out test set ($n\!=\!576$) only, to avoid
        train-set leakage.
\end{itemize}

Table~\ref{tab:erit-all} reports accuracy on the ERIT test set
($n=576$).  Across all model--modality combinations, fusion consistently
outperforms single-modality input.  LoRA fine-tuning provides substantial
gains: fusion MAcc rises from 60.9\,\% (LLaVA-NeXT 0-shot) to 70.0\,\%
(LoRA), a~+9.0\,pp improvement significant at $p < 0.001$ (McNemar).
Notably, LoRA text-only MAcc (64.1\,\%) exceeds both zero-shot
text-only baselines by a wide margin, indicating that supervised training
teaches the model to extract emotion cues from noisy Whisper transcripts.

\begin{table}[t]
\centering
\caption{Strict and multi-label accuracy~(\%) on the ERIT test set
($n{=}576$).  LoRA uses rank-8 adapters on the last 4 layers
(${\sim}$426K parameters).  McNemar's test: LoRA Fusion vs.\
zero-shot Fusion $p < 0.001$ for both models.}
\label{tab:erit-all}
\begin{tabular}{@{}llcccccc@{}}
\toprule
 & & \multicolumn{3}{c}{\textbf{Strict Acc (\%)}}
   & \multicolumn{3}{c}{\textbf{Multi-label Acc (\%)}} \\
\cmidrule(lr){3-5}\cmidrule(lr){6-8}
\textbf{Model} & \textbf{Regime} & I & T & I+T & I & T & I+T \\
\midrule
LLaVA-NeXT      & 0-shot & 30.4 & 16.8 & 33.3
                         & 49.8 & 29.7 & 60.9 \\
Qwen2-VL        & 0-shot & 33.3 & 20.5 & 35.6
                         & 46.4 & 34.6 & 51.9 \\
\midrule
LLaVA-NeXT LoRA & fine-tuned & 35.9 & 34.9 & \textbf{42.7}
                              & 58.0 & 64.1 & \textbf{70.0} \\
\bottomrule
\end{tabular}
\end{table}


% ─── 4.2 Robustness Analysis ─────────────────────────────────────────────

\subsection{Robustness Analysis}
\label{sec:robustness-results}

To probe model sensitivity to input quality and to disentangle the
contribution of each modality, we define five perturbation conditions
applied to the full evaluation set
(Table~\ref{tab:conditions}).  Each condition modifies one aspect of the
input while leaving the other modality intact, enabling controlled
measurement of robustness.

\begin{table}[t]
\centering
\caption{Perturbation conditions.  Each row describes a transformation
applied to all 2{,}379 samples before evaluation under the three modality
configurations.}
\label{tab:conditions}
\resizebox{\columnwidth}{!}{%
\begin{tabular}{llp{5.2cm}}
\toprule
\textbf{Condition} & \textbf{Modality} & \textbf{Transformation} \\
\midrule
Baseline         & ---   & Original images and transcripts \\
No-keyword       & Text  & Remove emotion-indicative words
                           (e.g., \emph{happy}, \emph{angry}, \emph{scared}) \\
Grayscale        & Image & Convert frames to grayscale
                           (removes colour-based cues such as skin flush) \\
Blurred faces    & Image & Apply Gaussian blur ($\sigma\!=\!23$)
                           to detected face regions
                           (probes reliance on facial morphology vs.\ context) \\
Text dropout     & Text  & Replace transcript with
                           ``No speech detected.'' \\
Image dropout    & Image & Replace frame with uniform gray image \\
\bottomrule
\end{tabular}}
\end{table}

\noindent
The combination of 6 conditions $\times$ 3 modality configurations yields
18 evaluation cells per model.  We exclude degenerate combinations where the
perturbation does not affect the active modality (e.g., grayscale under
text-only), resulting in 14 non-degenerate cells.

Table~\ref{tab:benchmark} reports LLaVA-NeXT zero-shot performance
across all 14 cells on the full dataset ($n=2{,}379$).

\begin{table}[t]
\centering
\caption{Robustness benchmark: LLaVA-NeXT zero-shot across 6
perturbation conditions and 3 modalities.  ``---'' denotes degenerate
cells (perturbation does not affect the active modality).
$\Delta$MAcc is the change from baseline.}
\label{tab:benchmark}
\resizebox{\columnwidth}{!}{%
\begin{tabular}{@{}l rrr rrr@{}}
\toprule
 & \multicolumn{3}{c}{\textbf{Multi-label Acc (\%)}}
 & \multicolumn{3}{c}{\textbf{$\Delta$MAcc (pp)}} \\
\cmidrule(lr){2-4}\cmidrule(lr){5-7}
\textbf{Condition} & I & T & I+T & I & T & I+T \\
\midrule
Baseline       & 59.2 & 37.4 & \textbf{69.5} &   --- &   --- &   --- \\
No-keyword     &  --- & 33.2 & 68.5           &   --- & $-$4.2 & $-$1.0 \\
Grayscale      & 56.6 &  --- & 68.6           & $-$2.6 &   --- & $-$0.9 \\
Blurred faces  & 24.4 &  --- & 55.0           & $-$34.8 & --- & $-$14.5 \\
Text dropout   &  --- &  3.4 & 54.2           &   --- & $-$34.0 & $-$15.3 \\
Image dropout  &  3.4 &  --- & 33.5           & $-$55.8 & --- & $-$36.0 \\
\bottomrule
\end{tabular}}
\end{table}

Three findings emerge.
\textbf{(i)~Fusion is resilient to mild perturbations:}
removing emotion keywords or converting to grayscale degrades fusion MAcc
by only 0.9--1.0\,pp, confirming that fusion draws on complementary cues.
\textbf{(ii)~Facial morphology is critical:}
blurring faces reduces image-only MAcc from 59.2\,\% to 24.4\,\%
($-$34.8\,pp), the largest single-modality drop, yet fusion recovers
to 55.0\,\%, demonstrating that text partially compensates for degraded
visual input.
\textbf{(iii)~Vision dominates:}
image dropout damages fusion far more than text dropout
($-$36.0 vs.\ $-$15.3\,pp), confirming that the visual modality carries
the stronger emotion signal.


% ─── 4.3 Stratified Analysis ─────────────────────────────────────────────

\subsection{Stratified Analysis}
\label{sec:stratified-results}

Beyond aggregate metrics, we report performance stratified along
diagnostically informative dimensions, computed post~hoc on the baseline
predictions without additional model inference.

\paragraph{Label agreement.}
Samples are split into \emph{high-agreement} (single annotator label;
$n\!=\!596$, 25.1\,\%) and \emph{low-agreement} subsets (two or more
valid labels; $n\!=\!1{,}783$, 74.9\,\%).

\paragraph{Modality conflict.}
We identify \emph{conflict samples} where the image-only and text-only
predictions disagree and further classify them as \emph{mild} (related
emotions) or \emph{opposite} (e.g., happiness vs.\ sadness).  This
subset isolates cases where fusion must actively reconcile contradictory
signals, making it the most informative slice for evaluating integration
quality.

\medskip
Table~\ref{tab:stratified} reports baseline fusion performance across
these slices.

\begin{table}[t]
\centering
\caption{Stratified analysis (baseline, LLaVA-NeXT 0-shot).
MAcc = multi-label accuracy.}
\label{tab:stratified}
\resizebox{\columnwidth}{!}{%
\begin{tabular}{@{}ll rrr r@{}}
\toprule
\textbf{Slice} & \textbf{Subset}
  & \textbf{I} & \textbf{T} & \textbf{I+T} & \textbf{$N$} \\
\midrule
\multicolumn{6}{l}{\textit{Label agreement (MAcc \%)}} \\
  & High (single label) & 37.6 & 23.3 & 36.7 & 596 \\
  & Low (multi-label)   & 66.5 & 42.1 & \textbf{80.5} & 1{,}783 \\
\midrule
\multicolumn{6}{l}{\textit{Modality conflict (MAcc \%)}} \\
  & Agree    & 54.3 & 54.3 & 65.8 & 622 \\
  & Mild     & 73.3 & 47.5 & \textbf{81.4} & 813 \\
  & Opposite & 50.3 & 17.5 & 61.8 & 944 \\
\bottomrule
\end{tabular}}
\end{table}

\paragraph{Label agreement results.}
Fusion MAcc on low-agreement (multi-label) samples reaches 80.5\,\%,
compared to 36.7\,\% on high-agreement samples.  This apparent paradox
arises because multi-label samples have a larger correct-answer set,
but it also shows that fusion reliably selects a plausible emotion
from the annotators' set---precisely the setting where integration
of visual and textual cues is most informative.

\paragraph{Modality conflict results.}
Of the 2{,}328 unique samples, 73.4\,\% exhibit a conflict between
image-only and text-only predictions; among those, 922 involve
\emph{opposite} emotions (e.g., image predicts surprise while text
predicts neutral; $n=631$, the dominant conflict pair).
Fusion recovers to 61.8\,\% MAcc on opposite-conflict samples despite
text-only accuracy collapsing to 17.5\,\%, indicating that the model
appropriately down-weights the weaker modality.
On mild-conflict samples, fusion achieves its highest MAcc (81.4\,\%),
suggesting that moderate cross-modal disagreement provides a useful
auxiliary signal.

\paragraph{Conflict pair analysis.}
The most frequent conflict pair is \emph{surprise}~(image) vs.\
\emph{neutral}~(text), accounting for 631 of 1{,}757 conflict samples
(35.9\,\%).  This aligns with the known difficulty of distinguishing
surprise from neutral in older adult faces, where raised eyebrows may
reflect attentiveness rather than surprise~\cite{ma2019elderreact}.


% ─── 4.4 Cross-Benchmark Comparison ──────────────────────────────────────

\subsection{Cross-Benchmark Comparison}
\label{sec:cross-benchmark}

To contextualise ERIT, Table~\ref{tab:fusion-only} compares fusion
performance with previously reported GPT-4V results on standard
multimodal benchmarks.  While GPT-4V achieves up to 80.4\,\% on CMU-MOSI
and 65.4\,\% on MER-MULTI, its ERIT fusion accuracy is 42.3\,\%.  This gap
reflects ERIT's controlled single-frame evaluation (no temporal
redundancy), the challenging domain of older adult expression, and the
high label ambiguity inherent in the dataset.
